% =======================================================================================
\section{The Conditioning Input}
\label{app:conditioning}
% =======================================================================================

The network sees the data of a window through three operations: whitening, a
Wilson--Daubechies--Meyer (WDM) transform, and a compression of the dynamic range.

\paragraph{Whitening.}
Each bin and channel is divided by its noise amplitude,
$\tilde d^{\,\rm w}_{k,c} = \tilde d_{k,c}/\sqrt{S_c(f_k)\Tobs/2}$, so that under noise alone
$\tilde d^{\,\rm w}_{k,c}\sim\CN(0,1)$: real and imaginary parts have variance $1/2$, and the
signal enters in units of its SNR per bin.

\paragraph{WDM transform.}
The WDM transform \citep{necula2012transient,cornish2020time,vajpeyi2026explicit} tiles the
time--frequency plane with $N_t$ time rows and $N_f$ frequency channels. From the one-sided
spectrum, channel $m$ draws on the Fourier bins within $N_t/2$ of $mN_t/2$:
\begin{equation}
  x_{n,m} = \sum_{l=-N_t/2}^{N_t/2-1}\tilde d^{\,\rm w}_{mN_t/2+l}\,\Phi_l\,e^{2\pi\ii nl/N_t},
  \qquad
  w_{n,m} = \sqrt2\,(-1)^{nm}\,\mathrm{Re}\bigl[C^*_{n,m}\,x_{n,m}\bigr],
  \qquad
  C_{n,m} = \begin{cases}1 & n+m\ \text{even}\\ \ii & n+m\ \text{odd},\end{cases}
  \label{eq:wdm}
\end{equation}
with $\Phi$ a Meyer-type window with a flat top over $|2l/N_t|\le a$ and a cosine taper of
width $b = 1-2a$, here $a = 1/3$. We compute only the channels that cover the window: $K$
channels starting at an even offset, as one reshape and a small matrix product, so that the
compiled graph does not grow with $K$ (a channel-by-channel loop compiles to millions of
operations on B).

We use the minimum $N_t = 2$ time rows, and in this case the transform simplifies exactly. The
window evaluates to $\Phi = (1, 0)$, so channel $m$ sees only bin $m$, and \cref{eq:wdm}
reduces to
\begin{equation}
  (w_{0,m},\,w_{1,m}) =
  \begin{cases}
    \sqrt2\,\bigl(\mathrm{Re}\,\tilde d^{\,\rm w}_m,\ \ \mathrm{Im}\,\tilde d^{\,\rm w}_m\bigr) & m\ \text{even},\\
    \sqrt2\,\bigl(\mathrm{Im}\,\tilde d^{\,\rm w}_m,\ -\mathrm{Re}\,\tilde d^{\,\rm w}_m\bigr) & m\ \text{odd},
  \end{cases}
  \label{eq:wdm-nt2}
\end{equation}
which we verified numerically against the implementation. The conditioning image is therefore
a lossless, real re-encoding of the whitened complex spectrum, two rows of unit-variance
pixels per TDI channel, rather than a time--frequency map. More time rows would trade
frequency resolution for time resolution, which the slowly evolving signals of this problem
do not need inside a window.
\draftnote{With $N_t=2$ the WDM step is a re-encoding of Re/Im. Present it that way in the
main text, or justify the WDM framing (e.g.\ a path to longer windows or other sources).}

\paragraph{Compression.}
A source puts $\sum_{k,c}|\tilde h^{\rm w}_{k,c}|^2 = \snr^2/2$ into a few bins, so the loudest
reach whitened amplitudes of order $10^3$ per bin, against unit-variance noise.
The input is $y = \operatorname{arcsinh}(w)$, linear in the noise and logarithmic for loud
signals, arranged as an image of $2\times K\times3$ pixels: time rows, frequency channels, and
the $A$, $E$, $T$ channels (\cref{fig:window}b). The same operations applied to the noiseless
signal give the target $y_{\rm clean}$ of the reconstruction head (\cref{eq:aux-x}).

\paragraph{What the image does not carry.}
Whitening and windowing remove the absolute frequency, which still matters: the Doppler width
scales with $f$, the chirp with $f_0^{11/3}$, and the amplitude window with $S_n(\fw)$. The
window centre $\fw$ is therefore given to the network separately (\cref{app:architecture}).

\begin{figure}[h]
  \centering
  \includegraphics[width=\textwidth]{figures/window.pdf}
  \caption{One XS window, from data to posterior (eval injection q0.20, final model). (a)
  Whitened power per bin, summed over $A$ and $E$, with the noiseless signal; vertical lines
  mark the true frequencies and their SNRs, and the shaded bins are the guard bands. (b) The
  conditioning input: $\operatorname{arcsinh}$ of the two-row WDM coefficients
  (\cref{eq:wdm-nt2}) for the three TDI channels. (c) The flow's $f_0$ posterior for each
  source, relabelled onto it (\cref{app:evaluation-matching}), against the Fisher forecast
  (each histogram scaled to a unit peak). The three detectable sources are found, with
  posteriors several times wider than the Fisher forecast. The SNR-5.3 source is not found,
  and its posterior does not spread over the window as the prior would: with exactly four
  sources per window and amplitudes of SNR $\gtrsim7$, a fourth source the data do not support
  can only sit where it is masked by the loud ones, which is where the flow puts it.}
  \label{fig:window}
\end{figure}
